\documentclass[11pt]{article}
\usepackage[utf8]{inputenc}
\usepackage{amsmath,amssymb,amsthm}
\usepackage[margin=1in]{geometry}
\usepackage{hyperref}
\usepackage{xcolor}
\usepackage{enumitem}
\usepackage{newunicodechar}
\newunicodechar{ℝ}{\ensuremath{\mathbb{R}}}
\newunicodechar{ℕ}{\ensuremath{\mathbb{N}}}
\newunicodechar{α}{\ensuremath{\alpha}}
\newunicodechar{γ}{\ensuremath{\gamma}}
\newunicodechar{λ}{\ensuremath{\lambda}}
\newunicodechar{μ}{\ensuremath{\mu}}
\newunicodechar{π}{\ensuremath{\pi}}

\title{Tide retrospective:\\\emph{J-function hypothesis weakening (L3)}}
\author{Daniel Murfet}
\date{7 May 2026}

\begin{document}
\maketitle

\begin{abstract}
This retrospective records a 152-line net deletion plus an
architectural unification of the J-function family. The N1 tide
(generic-$k$ J-function abstraction, merged earlier today) carried a
\texttt{Continuous prior} hypothesis stronger than the
mathematically-minimal one identified at the original E4
deliberation. This tide weakens the hypothesis to
\texttt{AEMeasurable prior volume + ContinuousAt prior 0 + global
bound} and, while the seabed is fresh, also refactors the sextic
specialisation\footnote{In \texttt{Laplace/OneD/JFunctionSextic.lean}
(L1, 2026-05-07).} to invoke the parametric form rather than carry
its own independent proof. The result: the parametric file gains
$\sim$32 lines of \texttt{quasiMeasurePreserving} plumbing; the
sextic file shrinks from 280 to 96 lines as its proofs become
five-line specialisations. Both Claude and GPT-5.5~Pro voted
Candidate B (weaken + refactor) on architectural grounds. Single
session, build green after one Mathlib-API-name correction
(\texttt{Measure.quasiMeasurePreserving\_smul} takes the measure as
an explicit argument, not implicit). Zero \texttt{sorry}, zero
\texttt{axiom}, zero \texttt{native\_decide}.
\end{abstract}

\section{Setting}

The grammar-paper precursor sequence has three closed Lean tides:
E4 (quartic J-function), L1 (sextic J-function), and N1 (generic-$k$
parametric form). All three statements take a continuous, bounded
prior $\pi : \mathbb R \to \mathbb R$. But the deliberation behind
E4 had GPT-5.5~Pro identify the mathematically-minimal hypothesis
class as
$$
  \pi \in \texttt{AEMeasurable} \;\wedge\;
  \pi \in \texttt{ContinuousAt}\,0 \;\wedge\;
  \exists M.\, \forall x.\, |\pi(x)| \le M.
$$
The stronger \texttt{Continuous prior} was taken in E4 / L1 / N1
because it was easier in Lean: composition with the scaling map
$u \mapsto u/t^\alpha$ at the AE-strong-measurability site is
trivial when $\pi$ is continuous.

L1's retrospective explicitly flagged the weakening as a follow-up,
estimating $\sim$50 lines of \texttt{comp\_quasiMeasurePreserving}
plumbing. N1's retrospective re-flagged it: \emph{``weakening should
land in N1 (\texttt{kth\_*}) only --- the specialisations inherit.''}

This tide takes both halves:

\begin{enumerate}
\item Weaken the parametric form's hypothesis class.
\item Refactor the sextic specialisation to invoke the parametric
  form (so it inherits the weaker hypothesis).
\end{enumerate}

(The quartic specialisation from E4 was never merged to \texttt{main}
--- it lived on the unmerged sibling branch and was superseded by N1
--- so the refactor scope is just sextic.)

\section{The thing formalised}

\begin{quote}\itshape
For the single-monomial potential $K(w) = w^{2k}/(2k)!$ with $k \ge 1$
and a globally bounded prior $\pi : \mathbb R \to \mathbb R$ that is
\texttt{AEMeasurable} and \texttt{ContinuousAt} 0,
\[
  \lim_{t \to \infty} t^{1/(2k)} \int_{\mathbb R}
    e^{-t\,w^{2k}/(2k)!}\,(\pi(w) - \pi(0))\,dw \;=\; 0,
\]
\[
  \lim_{t \to \infty} t^{1/(2k)} \int_{\mathbb R}
    e^{-t\,w^{2k}/(2k)!}\,\pi(w)\,dw \;=\;
    \pi(0) \cdot \tfrac{1}{k} \cdot ((2k)!)^{1/(2k)} \cdot \Gamma(1/(2k)).
\]
\end{quote}

\noindent The two parametric Lean signatures (now with the weakened
hypothesis class):
\begin{verbatim}
theorem kth_jfunction_centered_tendsto_zero
    {k : ℕ} (hk : 1 ≤ k)
    {prior : ℝ → ℝ}
    (hprior_meas : AEMeasurable prior volume)
    (hprior_cont0 : ContinuousAt prior 0)
    {M : ℝ} (hprior_bd : ∀ x, |prior x| ≤ M) :
    Tendsto ... atTop (𝓝 0)

theorem kth_jfunction_asymptotic
    {k : ℕ} (hk : 1 ≤ k)
    {prior : ℝ → ℝ}
    (hprior_meas : AEMeasurable prior volume)
    (hprior_cont0 : ContinuousAt prior 0)
    {M : ℝ} (hprior_bd : ∀ x, |prior x| ≤ M) :
    Tendsto ... atTop (𝓝 (prior 0 * ...))
\end{verbatim}

The two sextic theorems shrink from $\sim$80 lines each to $\sim$10
lines each, becoming thin specialisations at $k = 3$ with three
constant-bridging private lemmas
($(2 \cdot 3 : \mathbb N) = 6$ definitionally,
$(6! : \mathbb R) = 720$ via \texttt{decide} and
$1/((2 \cdot 3 : \mathbb N) : \mathbb R) = 1/6$ via \texttt{norm\_num}).

\section{Proof strategy}

The structural argument is unchanged from N1: substitute-then-DCT
with envelope $2M \cdot e^{-u^{2k}/(2k)!}$ and pointwise limit zero.
Only two sites in the proof use the old \texttt{Continuous prior}
hypothesis, both inside the centred theorem:

\paragraph{Site 1: AE strong measurability.} The DCT requires
strong AE measurability of the integrand
$F(t, u) := e^{-u^{2k}/(2k)!} \cdot (\pi(u/t^\alpha) - \pi(0))$
with respect to $t$. The old proof used the continuous-composition
upgrade. The weakened version routes through Mathlib's
quasi-measure-preserving composition for the scaling map
$u \mapsto (t^\alpha)^{-1} \cdot u$, which is canned for Haar-measure
spaces (and $\texttt{volume}$ on $\mathbb R$ qualifies).\footnote{Named
\texttt{AEMeasurable.comp\_quasiMeasurePreserving} with
\texttt{Measure.quasiMeasurePreserving\_smul}.} The reshape
$u / t^\alpha = (t^\alpha)^{-1} \cdot u$ takes the division into a
scalar action, making the canned scaling lemma applicable. Then
\texttt{AEMeasurable.aestronglyMeasurable} upgrades.

The eventual positivity of $t$ matters here:
$\texttt{Real.rpow}$ on a negative base does not give the expected
$t^\alpha \ne 0$ from $t \ne 0$, so the scaling factor's
nonzero-ness needs $0 < t$. The proof was already inside an
$\texttt{eventually\_gt\_atTop}$ filter; the fix is local.

\paragraph{Site 2: pointwise limit at $u$.} The DCT requires
$F(t, u) \to 0$ as $t \to \infty$ for a.e.\ $u$. The argument
$\pi(u/t^\alpha) \to \pi(0)$ comes from
$u/t^\alpha \to 0$ (\texttt{Tendsto.div\_atTop}) composed with
\texttt{Continuous.tendsto 0} (old) or
\texttt{ContinuousAt.tendsto} (new). One-line replacement.

\paragraph{The asymptotic theorem.} Inherits the new hypothesis bundle
without further proof changes (it's a corollary of the centred
theorem).

\paragraph{The sextic specialisation.} Each theorem reduces to a
five-line call: instantiate \texttt{kth\_*} at $k = 3$, then bridge
$(2 \cdot 3 : \mathbb N)$, $(2 \cdot 3)!$, and $1/((2 \cdot 3 : \mathbb N) : \mathbb R)$
to their numeric forms via three \texttt{simp\_rw}-firing lemmas.

\section{Roadblocks and resolutions}

One Mathlib-API-name correction during the build, two tactical
adjustments, and the standard \texttt{simp\_rw} pattern for the
specialisation bridges.

\paragraph{Mathlib lemma signature.}
The canned scaling-quasi-MP lemma takes the measure $\mu$ as an
\emph{explicit} positional argument (it lives inside a section with
an explicit \texttt{variable (\mu : Measure E)}), not implicit.
First-attempt call without passing $\mu$ gave the diagnostic
\emph{``argument has type \texttt{$\dots$ \neq 0} but is expected to
have type \texttt{Measure ?m}''} --- the type mismatch was the
clue. Fix: pass $\mu = \texttt{volume}$ as the first argument.
Five-second adjustment.

\paragraph{Sextic specialisation: \texttt{convert} versus
\texttt{simp\_rw}.} First attempt at the sextic refactor used
\texttt{convert h using 2} to align the parametric form's
\texttt{Tendsto} statement with the sextic-specific form. This
left several minor goals that needed independent constant-bridging
proofs, ending in \emph{``No goals to be solved''} as the convert
closed cleanly but the trailing tactics fired anyway. Cleaner
approach: do all constant-bridging via \texttt{simp\_rw at h} on the
parametric hypothesis, then close with \texttt{exact h}. Three
\texttt{simp\_rw}-shaped private lemmas
(for the $\mathbb N$-power, the factorial, and the rational inverse)
suffice; the underlying integrand and limit forms align by
$\texttt{rfl}$ once all three are bridged.

\paragraph{No GPT thrashing-rescue.} The deliberation at Step~2 named
the right Mathlib lemma (\texttt{quasiMeasurePreserving\_smul}) and
the right shape; the Step~3 frictions were name-and-form bookkeeping,
not strategy errors.

\section{What was Mathlib, what was new}

\paragraph{Mathlib.} The scaling quasi-MP fact for Haar measure on
finite-dim $\mathbb R$ vector spaces, the AE-measurability transport
under quasi-MP composition, the upgrade from AE measurability to AE
strong measurability for $\mathbb R$-valued functions, and the
standard \texttt{ContinuousAt.tendsto} fact. All familiar; no new
Mathlib API surfaced.

\paragraph{Laplace seabed (existing).} N1's parametric \texttt{kth\_*}
theorems (now weakened in this tide) and the supporting private
helpers (\texttt{kth\_jfunction\_centered\_subst},
\texttt{kth\_dominator\_integrable}, etc.) --- unchanged. L1's
sextic-specific helpers
(\texttt{sextic\_jfunction\_centered\_subst},
\texttt{sextic\_dominator\_integrable}) ---
\emph{deleted} as part of the refactor; the parametric versions
subsume them.

\paragraph{New.} Three constant-bridging private lemmas in the
refactored sextic file (numeric specialisations of $k = 3$). The
parametric file gains $\sim$32 lines of inline
quasi-MP plumbing inside
\texttt{kth\_jfunction\_centered\_tendsto\_zero} (helper not factored
out, since only one use site; if a third $k$-specialisation tide
materialises, extracting an
\texttt{aemeasurable\_comp\_div\_const} helper would pay off).

\section{Lessons}

\begin{itemize}
\item \emph{The minimal-hypothesis-class deliberation pays off late,
  not early.} GPT-5.5~Pro identified the
  \texttt{AEMeasurable + ContinuousAt 0 + bound} class at E4's
  Step~2 deliberation; E4, L1, and N1 all chose the stronger
  \texttt{Continuous prior} for Lean ergonomics. This tide redeemed
  the deferred work; the cost was small ($\sim$32 lines of
  parametric-side plumbing) but the architectural payoff
  ($-184$ lines on the sextic side, sextic now inherits future
  parametric improvements automatically) was much larger.
  Lesson: the minimal hypothesis class is worth landing as a
  follow-up tide once the seabed is fresh, even when the original
  tide takes the easier-in-Lean form. Don't make a religion of it
  during the original tide.
\item \emph{\texttt{simp\_rw at h} beats \texttt{convert} for
  specialisation bridging.} The sextic specialisation has the same
  underlying mathematical content as $k = 3$ in the parametric form;
  the difference is purely numeric ($(2 \cdot 3 : \mathbb N) = 6$ etc.).
  Several rounds of \texttt{simp\_rw at h} followed by
  \texttt{exact h} is more robust than \texttt{convert} for this
  pattern, because \texttt{convert} leaves goals at depth-mismatch
  and forces extra work; \texttt{simp\_rw at h} drives the rewrites
  in the right direction without leaving leftovers.
\item \emph{Mathlib lemma signatures vary in implicit/explicit
  bindings.} \texttt{Measure.quasiMeasurePreserving\_smul} takes the
  measure as an explicit positional argument; many other
  \texttt{quasiMeasurePreserving\_*} lemmas take it implicit.
  This is normal Mathlib --- the explicit-binding pattern reflects
  the section's \texttt{variable} declarations --- but it's a
  one-attempt-then-fix cost. Worth a CLAUDE.md note for future
  tides.
\end{itemize}

\section{Follow-ups}

\begin{itemize}
\item \emph{Promote \texttt{aemeasurable\_comp\_div\_const}.} If a
  third hypothesis-weakening tide arrives (e.g., for a tide outside
  the J-function family that uses the same $u/c$ scaling pattern),
  factor the inline plumbing of this tide into a public helper. Two
  use sites is currently below the refactor threshold; three would
  be over.
\item \emph{Weaken further to AE bound.} The current
  \texttt{$\forall$ x.~|prior x| $\le$ M} could be relaxed to
  \texttt{$\forall^\mu$ x.~|prior x| $\le$ M} (essentially bounded)
  with minor proof changes. GPT explicitly said \emph{``not this
  tide''}; mark as a refinement candidate for a future
  hypothesis-weakening tide.
\item \emph{\texttt{\textbackslash leanref}-tag the grammar paper
  with the weaker hypothesis class.} Once the grammar paper has TeX
  statements, both the centred and asymptotic theorems should be
  pinned to this commit and tagged with the weakened-hypothesis
  form.
\item \emph{Promote \texttt{kth\_jfunction\_centered\_subst} and
  \texttt{kth\_dominator\_integrable} to \texttt{lean-common}.}
  Now that the sextic file no longer has its own copies, these are
  the only Laplace-specific J-function helpers; they're potentially
  useful for any monomial-Gibbs Laplace-method tide. Defer until a
  consumer outside laplace materialises.
\end{itemize}

\section*{Acknowledgements}

GPT-5.5~Pro's deliberation contributed three load-bearing pieces of
guidance. First, the recommendation to use \texttt{AEMeasurable}
rather than \texttt{AEStronglyMeasurable} as the public hypothesis
(\emph{``\texttt{AEStronglyMeasurable} is proof-shaped rather than
theorem-shaped; recover it internally via
\texttt{aestronglyMeasurable}''}). Second, the architectural vote for
Candidate B (refactor sextic), with the practical suggestion of
two-commits-one-PR packaging (collapsed here for git simplicity).
Third, the Mathlib-API-name candidates and the eventual-positivity
caveat for \texttt{Real.rpow}. Full consult preserved at the GPT v1
file.\footnote{Preserved at
\texttt{projects/primer/tide-log/gpt55\_tide\_jfunction\_hypothesis\_weakening\_v1.md}.}

This Tide branched off the post-O1-merge \texttt{main} (commit
\texttt{dc9d260}). No in-flight tides interact with this work; the
parametric file's only consumer is the sextic file, which is also
modified here.

\end{document}
